% The manual of the algomotion LaTeX package. Its figures are drawn by the
% package itself, so it compiles like any document that uses it:
%
%   lualatex -shell-escape algomotion-doc.tex   (twice, for the contents)
%
% with the algomotion command on the PATH. In the algomotion repository,
% `npm run build:doc -w @algomotion/latex` builds it with the command of
% the checkout.
\documentclass[11pt]{article}
\usepackage[margin=2.5cm]{geometry}
\usepackage{algomotion}
\usepackage{booktabs}
\usepackage{tabularx}
\usepackage{listings}
\usepackage{xcolor}
\usepackage[hidelinks]{hyperref}

\lstset{
  basicstyle=\ttfamily\small,
  columns=fullflexible,
  keepspaces,
  breaklines,
  frame=leftline,
  framerule=1.5pt,
  rulecolor=\color{black!25},
  xleftmargin=1.5em,
  aboveskip=0.8\baselineskip,
  belowskip=0.8\baselineskip,
}
\newcommand\pkg[1]{\textsf{#1}}
\newcommand\cs[1]{\texttt{\textbackslash#1}}
\newcommand\key[1]{\texttt{#1}}
\makeatletter
\edef\algomotionversion{\csname ver@algomotion.sty\endcsname}
\makeatother

\title{The \pkg{algomotion} package\\[0.3em]
  \large Step-by-step figures of algorithms and data structures}
\author{Eryc123Y\\
  \url{https://github.com/Eryc123Y/algorithm-visualizer}}
\date{\algomotionversion}

\begin{document}
\maketitle

\begin{abstract}
  Name an algorithm, give it your input and pick a step, such as the moment
  Dijkstra's algorithm improves a distance on your graph, and \pkg{algomotion}
  draws that step as a TikZ figure, in your document's fonts. It can also draw
  a document that you wrote, or that an AI wrote for you, in the algomotion
  format. In Beamer, a figure becomes overlays that step through the run.
  Every figure in this manual is drawn by the package.
\end{abstract}

\tableofcontents

\section{Introduction}

\noindent This line draws figure~\ref{fig:dijkstra}:

\begin{lstlisting}
\algomotion[step=improve#3]{dijkstra}{A-B 4, B-C 1, A-C 7, C-D 2}
\end{lstlisting}

\begin{figure}[htbp]
  \centering
  \algomotion[step=improve\#3, layout=wide, width=0.75\linewidth]{dijkstra}{A-B 4, B-C 1, A-C 7, C-D 2}
  \caption{Dijkstra's algorithm from $A$, the third time a distance improves:
    the edge from $B$ lowers the distance of $C$ from $7$ to $5$. The
    priority queue is drawn beside the graph.}
  \label{fig:dijkstra}
\end{figure}

\noindent The graph is an edge list, and Dijkstra's algorithm starts at its
first node, $A$. The figure itself uses \texttt{layout=wide}, which puts the
priority queue beside the graph.

The figures come from algomotion, a library and command line that run the
algorithms of a registry, about 130 of them from sorting to graphs, dynamic
programming and data structures, and record every step. The same figures are
drawn by its web player, its Typst package and its SVG export, with the same
shapes, colours and positions. Here they are drawn with TikZ, so their text
is set in the document's fonts and they scale with the page.

The package runs the \texttt{algomotion} command through shell escape, as
\pkg{minted} runs Pygments. It keeps every figure in a cache directory, so a
figure that did not change is never made again, and cached figures need no
shell escape at all.

\section{Installation}

\begin{itemize}
  \item \textbf{The command.} Install \texttt{algomotion} with npm,
    \texttt{npm install -g @algomotion/mcp}, or take the standalone
    executable for your platform from the project's GitHub releases, a
    program of about 3\,MB that needs no Node.js. Put it on the
    \texttt{PATH}, or name it with the package option \key{command}.
  \item \textbf{The package.} \texttt{algomotion.sty} comes with
    \TeX\ distributions that include it from CTAN; otherwise put it next to
    your document, or in your local \texttt{texmf} tree under
    \nolinkurl{tex/latex/algomotion/}.
  \item \textbf{Requirements.} \LaTeX\ 2024-06-01 or later, as in \TeX\ Live
    2024, and TikZ. pdf\LaTeX, Lua\LaTeX\ and Xe\LaTeX\ work alike;
    \texttt{latex} with \texttt{dvips} works too, but PostScript has no
    transparency, so the shadows of the light and dark themes come out solid.
  \item \textbf{Shell escape} is needed the first time a figure is made:
\end{itemize}

\begin{lstlisting}
latexmk -pdf -shell-escape notes.tex     % or -pdflua, -pdfxe
\end{lstlisting}

\section{Commands}

\begin{center}
  \begin{tabularx}{\linewidth}{@{}lX@{}}
    \toprule
    Command & Draws \\
    \midrule
    \cs{algomotion}\texttt{[keys]\{id\}\{input\}} & a step of a registry
      algorithm, run on the input (section~\ref{sec:input}) \\
    \cs{algomotionfile}\texttt{[keys]\{doc.json\}} & a step of an algomotion
      document in a JSON file \\
    \cs{algomotionfile}\texttt{[algorithm=id, keys]\{in.json\}} & a step of
      an algorithm whose input is in a file \\
    \cs{algomotioninput}\texttt{[keys]\{figure.tex\}} & a figure or series
      made by \texttt{algomotion export --format tikz}, without shell escape \\
    \cs{algomotionset}\texttt{\{keys\}} & sets keys for the figures that
      follow, in the current group \\
    \bottomrule
  \end{tabularx}
\end{center}

\noindent \texttt{algomotion recipes} lists the algorithm ids and the inputs
they take, and \texttt{algomotion guide} explains documents.

\section{Keys}

\begin{center}
  \small
  \begin{tabularx}{\linewidth}{@{}llX@{}}
    \toprule
    Key & Default & Meaning \\
    \midrule
    \key{step} & \texttt{-1} & The frame to draw (section~\ref{sec:steps}). \\
    \key{steps} & none & A series of frames, one per slide in Beamer
      (section~\ref{sec:slides}). \\
    \key{start} & \texttt{+} & The slide of a series' first frame. \\
    \key{theme} & \texttt{light} & A built-in theme
      (section~\ref{sec:themes}). \\
    \key{layout} & \texttt{narrow} & Where structures beside the main one go,
      such as a queue: below it (\texttt{narrow}) or beside it
      (\texttt{wide}); Beamer's default is \texttt{wide}. \\
    \key{lang} & \texttt{en} & The language of an algorithm's run and of the
      figure's own words: \texttt{en} or \texttt{zh} (Simplified Chinese,
      in a document that sets Chinese, such as one with \pkg{ctex}). \\
    \key{width}, \key{height} & natural & The figure's size; with both, it
      fits inside both and keeps its shape. \\
    \key{scale} & none & A factor on the natural size. \\
    \key{value font} & the theme's & The font of values in cells and nodes,
      indices, edge labels and pointers. \\
    \key{label font} & the theme's & The font of titles and labels such as
      ``top''. \\
    \key{algorithm} & none & With \cs{algomotionfile}: the file is this
      algorithm's input, not a document. \\
    \bottomrule
  \end{tabularx}
\end{center}

An algorithm's figure also draws what it keeps besides its main structure,
such as Dijkstra's priority queue or a call stack. \texttt{layout=narrow} puts
those structures below the main one, which suits a column of notes;
\texttt{layout=wide} puts them beside it, which suits a slide.

At its natural size, a figure's values come out at 1.25 times the size of the
text around it, and its labels at three quarters of it. A figure is never
wider than the line, unless \key{width} or \key{height} says otherwise. Every
figure of one run of an algorithm has the same canvas, so figures of
different steps drawn at the same scale line up:

\begin{lstlisting}
\algomotion[step=swap#2, width=0.45\linewidth]{bubble-sort}{5 1 4 2 8}\hfill
\algomotion[step=-1, width=0.45\linewidth]{bubble-sort}{5 1 4 2 8}
\end{lstlisting}

\begin{figure}[htbp]
  \centering
  \algomotion[step=swap\#2, width=0.45\linewidth]{bubble-sort}{5 1 4 2 8}\hfill
  \algomotion[step=-1, width=0.45\linewidth]{bubble-sort}{5 1 4 2 8}
  \caption{Bubble sort's second swap and its last frame, at one width and so
    at one scale.}
\end{figure}

\noindent Package options, in \cs{usepackage}\texttt{[...]\{algomotion\}}:
\key{command=algomotion}, the command to run, with a path or arguments if
needed; and \key{cachedir=\_algomotion}, where the figures are kept.

\section{Steps}\label{sec:steps}

\key{step} takes a selector:
\begin{itemize}
  \item a frame number: \texttt{0} is the starting state, \texttt{-1} the
    last frame;
  \item a step name, counted from 1: \texttt{relax\#3} is the third step
    named \texttt{relax};
  \item a step inside a section: \texttt{insert-5/relax\#2}.
\end{itemize}
A selector that names several frames, such as a section or a range, is an
error with \key{step}, since a figure is one frame; \key{steps} takes a
series. If a step does not exist, the error lists the step names the run has.

In ordinary text, write \texttt{step=relax\#3}. Inside a macro definition, a
table or a Beamer frame without \texttt{[fragile]}, \TeX\ reads \texttt{\#}
as a parameter, so write \texttt{step=relax\textbackslash\#3} there. That
spelling works everywhere.

\section{Themes}\label{sec:themes}

Nine themes are built in: \texttt{light} (the default), \texttt{dark},
\texttt{print}, \texttt{textbook}, \texttt{textbook-dark}, \texttt{blocks},
\texttt{blocks-dark}, \texttt{answer-sheet} and \texttt{turbo-c}. Each one
chooses the document's family for its values and labels, which the keys
\key{value font} and \key{label font} change. The dark themes and Turbo~C
paint their own background. Figure~\ref{fig:themes} draws the step of
figure~\ref{fig:dijkstra} in six of them.

\newcommand\themed[1]{%
  \begin{minipage}[b]{0.31\linewidth}
    \centering
    \algomotion[step=improve\#3, theme=#1, layout=wide, width=\linewidth]{dijkstra}{A-B 4, B-C 1, A-C 7, C-D 2}\\[0.3em]
    \texttt{#1}
  \end{minipage}}
\begin{figure}[htbp]
  \centering
  \themed{print}\hfill\themed{textbook}\hfill\themed{blocks}\\[1.2em]
  \themed{answer-sheet}\hfill\themed{textbook-dark}\hfill\themed{turbo-c}
  \caption{One step in six themes.}
  \label{fig:themes}
\end{figure}

\section{Slides}\label{sec:slides}

In a Beamer frame, \key{steps} turns a figure into overlays: one slide per
frame of a series, all in one picture, so nothing moves from slide to slide.

\begin{lstlisting}
\begin{frame}{Dijkstra's algorithm, step by step}
  \centering
  % every frame
  \algomotion[steps]{dijkstra}{s-a 4, s-b 1, b-a 2, a-c 1, b-c 5, c-t 3}
\end{frame}

\begin{frame}{Settling $a$}
  The queue hands out $a$, at distance $3$. \pause

  \centering
  \algomotion[steps=extract\#3..improve\#5]{dijkstra}{s-a 4, s-b 1, b-a 2,
    a-c 1, b-c 5, c-t 3}
\end{frame}
\end{lstlisting}

\begin{itemize}
  \item \key{steps} takes the selectors of \key{step}, and series of them: a
    range such as \texttt{relax\textbackslash\#2..-1}, a section such as
    \texttt{insert-5}, or several in braces:
    \texttt{steps=\{0, relax\textbackslash\#2, -1\}}. Without a value, it
    means every frame. The frames come in timeline order, each once.
  \item A series starts where Beamer's incremental count stands, as
    \texttt{<+>} does, so it follows a \cs{pause}; \texttt{start=3} starts it
    on slide 3 instead. Its last frame stays until the frame ends; continue
    with \cs{onslide}\texttt{<+->} rather than \cs{pause}.
  \item In the \texttt{beamer} class, a figure fills the frame: as large as
    fits in the line's width and three quarters of the text height, keeping
    its shape. \key{width} and \key{height} replace their side of that box,
    and \key{scale} gives the natural size times \key{scale}.
  \item In handouts, in transparencies, in \texttt{beamerarticle} notes and
    outside Beamer, a series shows its last frame.
\end{itemize}

\section{The input}\label{sec:input}

The input of \cs{algomotion} is a notation, or JSON:

\begin{lstlisting}
\algomotion{insertion-sort}{5 2 4 6 1 3}                     % a list
\algomotion{dijkstra}{A-B 4, B-C 1, A-C 7, C-D 2}            % a graph
\algomotion{topological-sort}{shirt->tie, tie->jacket, socks} % a directed graph
\algomotion{inorder-tree-walk}{[4, 2, 6, 1, 3, null, 7]}     % a tree in level order
\algomotion{bst}{"from": "8 3 10 1 6", "ops": "insert 4, delete 3"}
\algomotion{dijkstra}{"graph": "A-B 4, B-C 1", "source": "C"}
\algomotion{kruskal}{}                                        % the first example
\end{lstlisting}

\begin{itemize}
  \item \textbf{Lists} are numbers separated by commas or spaces, optionally
    in brackets; a level order writes \texttt{null} for a gap.
  \item \textbf{Graphs} are edge lists: \texttt{-} joins an undirected edge
    and \texttt{->} a directed one, a number after a space is the weight, and
    a name alone is a node without edges.
  \item \textbf{Operations} of a data structure are written as exercises
    write them: \texttt{insert 4, delete 3}.
  \item Anything that starts with \texttt{\{}, \texttt{[} or \texttt{"} is
    JSON, and the outer braces of an object may be left out.
  \item Write \texttt{\%} as \texttt{\textbackslash\%}; \texttt{\textbackslash\#},
    \texttt{\textbackslash\{}, \texttt{\textbackslash\}}, \texttt{\textbackslash\&},
    \texttt{\textbackslash\$} and \texttt{\textbackslash\_} also stand for
    the characters. Line breaks and repeated spaces are fine.
\end{itemize}

\noindent Figure~\ref{fig:inputs} draws two of these inputs, and a dynamic
programming table.

\begin{figure}[htbp]
  \centering
  \begin{minipage}[c]{0.4\linewidth}
    \centering
    \algomotion[width=\linewidth]{bst}{"from": "8 3 10 1 6", "ops": "insert 4, delete 3"}
  \end{minipage}\hfill
  \begin{minipage}[c]{0.55\linewidth}
    \centering
    \algomotion[width=\linewidth]{lcs}{"first": "ABCB", "second": "BDCAB"}
  \end{minipage}
  \caption{Left, a binary search tree built from \texttt{8 3 10 1 6}, after
    inserting $4$ and deleting $3$. Right, the table of the longest common
    subsequence of \texttt{ABCB} and \texttt{BDCAB}.}
  \label{fig:inputs}
\end{figure}

\section{The cache}

Each figure is kept in \texttt{\_algomotion/} under the MD5 of everything
that changes its drawing: the input or the file's contents, the algorithm,
the step, the theme, the layout, the file format and the package version.

\begin{itemize}
  \item A figure found there is drawn without running anything, so later
    runs of a document cost only TikZ's time.
  \item Fonts and sizes are applied when the figure is drawn, so changing
    them makes no new figure.
  \item Share or commit \texttt{\_algomotion/} to compile the document where
    the command is missing or shell escape is off, such as on arXiv or for a
    co-author.
  \item The package deletes nothing. Delete \texttt{\_algomotion/} to clean
    up, and after updating the \texttt{algomotion} command.
\end{itemize}

\section{Errors}

When a figure cannot be made, \LaTeX\ reports an error, and a framed box
takes the figure's place with the whole message: an unknown algorithm or
theme, an input that the algorithm rejects, an invalid document, or a step
that names no frame. The package also says when shell escape is off and the
figure is not in the cache yet; when restricted shell escape, the \TeX\ Live
default, did not run \texttt{algomotion}, which is not in
\texttt{shell\_escape\_commands}; and when the command made no figure,
usually because it is not installed or not on the \texttt{PATH}.

\section{Security}

Shell escape lets a document run programs, so compile documents from others
with care, as with \pkg{minted}. The package passes the shell only checked
text: the command line holds the algorithm id, the step and the theme, each
checked against the characters it may have; the input is passed as
hexadecimal; and file names are quoted, and rejected if they contain
characters that shells read inside quotes.

\section{Licence and support}

The package is free software under the GNU General Public License, version 3
or later. The figures it draws are yours.

\begin{itemize}
  \item Report problems and ask questions at\\
    \url{https://github.com/Eryc123Y/algorithm-visualizer/issues}.
  \item The documentation of the whole project, with every algorithm and its
    figures, is at \url{https://algomotion.turing2234.com}.
\end{itemize}

\end{document}
